\chapter{The render pipeline end-to-end}
\label{ch:pipeline2}

\begin{aside}
You have seen each stage. This chapter walks one frame from
``user pressed a key'' to ``ANSI on the wire'' with all the
parts assembled.
\end{aside}

\section{The key press}

The user holds \code{j} in a list widget to scroll. The kernel
delivers the byte to \maya{}'s pty reader. The reader's
\code{poll} returns; the loop wakes.

\section{Parse}

The byte is fed to the input parser. The parser is a state
machine that handles escape sequences, modifier folding, and
bracketed paste. For a bare \code{j} byte it emits one event:

\begin{lstlisting}
KeyEvent { code = 'j', special = None }
\end{lstlisting}

\section{Dispatch}

The event reaches the dispatcher, which walks the focus stack.
The list widget has focus; its key handler matches \code{j}
to ``move selection down''.

\section{Mutation}

The handler calls \code{ctx.set(selection, sel + 1)}. This:

\begin{enumerate}[leftmargin=*]
  \item Updates the signal's stored value.
  \item Walks the dependency graph and marks every effect
    that read this signal as dirty.
  \item Schedules a frame.
\end{enumerate}

\section{Render}

The scheduler dequeues the frame:

\begin{itemize}[leftmargin=*]
  \item \textbf{Layout.} The element tree is walked; each
    element computes its size and position. The list's height
    is the parent height; the visible rows are the slice
    starting at scroll offset.
  \item \textbf{Paint.} The back buffer is cleared and the
    tree is painted depth-first. The list draws its rows;
    one row is highlighted (the new selection).
  \item \textbf{Diff.} The back and front buffers are
    compared. Two rows differ: the previously selected row
    (now unhighlighted) and the new one.
  \item \textbf{Emit.} ANSI for the two rows is generated and
    bundled into one string.
\end{itemize}

\section{Write}

The string goes to \code{write(1, ...)}. The kernel queues it
to the pty. The terminal emulator reads it from the pty
master, parses the ANSI, and updates its internal screen
buffer.

\section{Pixels}

The terminal emulator renders the screen buffer to pixels
using its font and palette. The window manager composites the
window. The GPU writes to the monitor. The user's eye sees
the change.

\section{Time budget}

For a key-press-to-screen latency:

\begin{itemize}[leftmargin=*]
  \item Kernel + pty: <1 ms.
  \item Parse: ~10 microseconds.
  \item Dispatch + mutation: ~10 microseconds.
  \item Layout: 0.5 ms typical.
  \item Paint: 1 ms typical.
  \item Diff: 0.5 ms typical.
  \item Emit + write: 0.5 ms.
  \item Terminal render: 5--15 ms depending on terminal.
\end{itemize}

Total: under 20 ms, well within a 60 Hz visual budget.

\section{What slows it down}

\begin{itemize}[leftmargin=*]
  \item Big trees: layout is O(N) but N can be thousands.
  \item Big screens: paint and diff are O(cells).
  \item Slow terminals: render dominates on terminals
    without GPU acceleration.
  \item Bad coalescing: handling one key at a time when many
    arrive in a burst.
\end{itemize}

The renderer fixes the first three; the loop's coalescing
fixes the fourth.

\section{Profiling}

\maya{} can dump per-stage timings to a trace file. Enable
with \code{MAYA\_TRACE=1}. The file is JSON; tools like
chrome://tracing or Perfetto can visualise it.

\begin{tryit}
Enable tracing and capture 10 seconds of a typing burst.
Look at the timeline. Layout and paint should both be sub-
millisecond; if one spikes, that is your bottleneck.
\end{tryit}

\section{What if it crashes mid-frame}

A crash during paint leaves the back buffer in an undefined
state. The atexit hook still runs: shutdown bundle restores
the terminal. The user sees their pre-crash shell.

A crash during emit (mid-write) is harmless; the kernel
delivers whatever bytes made it, and shutdown bytes are sent
after.

\section{Recap}

\begin{recap}
\begin{itemize}[leftmargin=*]
  \item Key to pixel in under 20 ms; under 5 in \maya{} land.
  \item Each stage is bounded and instrumentable.
  \item Coalescing keeps wakeups proportional to user
    actions, not byte arrivals.
\end{itemize}
\end{recap}

\section*{Workshop}

\begin{enumerate}[leftmargin=*]
  \item Enable tracing and identify the slowest stage in
    your app.
  \item Type a paragraph of text rapidly; confirm one
    frame per multi-key burst, not one per key.
  \item Measure end-to-end latency with a stopwatch app
    on your phone. Compare your terminal with another.
\end{enumerate}
